\documentclass[11pt,a4paper]{article}
\usepackage[utf8]{inputenc}
\usepackage[T1]{fontenc}
\usepackage{lmodern}
\usepackage[spanish,es-tabla]{babel}
\usepackage[margin=2.2cm]{geometry}
\usepackage{graphicx}
\usepackage{booktabs}
\usepackage{tabularx}
\usepackage{tabularray}
\usepackage{array}
\usepackage{titlesec}
\usepackage{enumitem}
\usepackage{microtype}
\usepackage{hyperref}

\hypersetup{
    colorlinks=true,
    linkcolor=black,
    urlcolor=blue,
    citecolor=black
}

\titleformat{\section}{\large\bfseries}{}{0em}{}[\titlerule]
\titleformat{\subsection}{\normalsize\bfseries}{}{0em}{}

\setlist[itemize]{noitemsep, topsep=2pt}
\setlist[enumerate]{noitemsep, topsep=2pt}

\begin{document}
\hypersetup{pageanchor=false}

% =========================================================
% PORTADA OFICIAL (PÁGINA INDEPENDIENTE)
% =========================================================
\begin{titlepage}
    \centering
    
    \includegraphics[width=0.28\textwidth]{logo.png} 
    \vspace{1.2cm}
    
    {\fontsize{13}{16}\selectfont \textbf{UNIVERSIDAD NACIONAL MAYOR DE SAN MARCOS}}\\[0.3em]
    {\normalsize Universidad del Perú. Decana de América}\\[1.5em]
    
    {\large \textbf{``Año de la Esperanza y el Fortalecimiento de la Democracia''}}\\[3.5em]
    
    {\Large \textbf{MÉTODOS CUANTITATIVOS DE ANÁLISIS POLÍTICO}}\\[2.5em]
    
    {\LARGE \textbf{\textit{SÍLABO}}}\\[4.5em]
    
    \begin{flushleft}
    \fontsize{11}{18}\selectfont
    \begin{tabular}{ll}
        \textbf{Denominación del programa:} & Doctorado en Ciencia Política \\
        \textbf{Profesor responsable:}      & Dr. José Manuel Magallanes Reyes, Ph.D. \\
        \textbf{Correo electrónico:}        & jmagallanesr@unmsm.edu.pe\\
    \end{tabular}
    \end{flushleft}
    
    \vfill
    
    {\LARGE \textbf{2026}}
    
\end{titlepage}

\hypersetup{pageanchor=true}
\setcounter{page}{1}

% =========================================================
% CUERPO DEL SÍLABO
% =========================================================

\section{I. INFORMACIÓN GENERAL}

\noindent\begin{tabularx}{\textwidth}{>{\bfseries}p{6.2cm}X}
Nombre de la asignatura: & Métodos Cuantitativos de Análisis Político \\
Código: & D4S10122 \\
Tipo de asignatura: & Obligatoria / Especialidad \\
Docente: & Dr. José Manuel Magallanes Reyes, Ph.D. \\
Programa: & Doctorado en Ciencia Política \\
Créditos: & 4.0 \\
Horas semanales: & 3 horas (Teórico-Prácticas / Taller) \\
Horas por módulo: & 48 horas totales (16 sesiones sabatinas $\times$ 3 horas) \\
Semestre académico: & 2026-II \\
Duración: & 16 semanas \\
Fecha de inicio: & 12 de Setiembre de 2026 \\
Fecha de finalización: & 26 de Diciembre de 2026 \\
Horario: & Sábados de 16:00 -- 18:15 (sesiones virtuales, por grupo, cada dos semanas) \\
\end{tabularx}

\section{II. FUNDAMENTOS DE LA ASIGNATURA}

\subsection{Sumilla}
Asignatura de naturaleza teórico-práctica y de carácter avanzado orientada a la investigación doctoral. El seminario introduce y profundiza en cinco rutas metodológicas de la ciencia política empírica contemporánea: (1) Regresión Discontinua en el Tiempo (RDiT), (2) Análisis de Duración y Riesgo (Event History Analysis -- EHA), (3) Métodos Configuracionales para Muestras Pequeñas (Fuzzy-set QCA), (4) Medición Computacional y Procesamiento de Lenguaje Natural (Text-as-Data), y (5) Teoría de Juegos y Modelado Formal de Decisiones Estratégicas. A través de una pedagogía de taller, el estudiante recorre una guía de autoestudio con notebooks reproducibles, elige una única ruta metodológica para su proyecto de tesis apoyándose en una lectura de referencia asociada a esa ruta, y desarrolla un artículo de investigación empírica original (\textit{working paper}) con código reproducible, cuya nota se construye progresivamente a partir del avance evaluado en las asesorías sabatinas.

\subsection{Perfil de Egreso}
El egresado del Doctorado en Ciencia Política contará con el siguiente perfil:
\begin{itemize}
    \item Generará conocimientos innovadores en Ciencia Política para resolver desafíos políticos y de gobernanza contemporáneos, con rigor metodológico y perspectiva crítica.
    \item Analizará problemáticas políticas complejas para proponer soluciones basadas en evidencia, aplicando métodos de investigación avanzados y pensamiento sistémico.
    \item Desarrollará marcos teóricos y metodológicos en Ciencia Política para contribuir a la evolución de la disciplina, con originalidad y fundamentación epistemológica sólida.
    \item Liderará proyectos de investigación y asesoría en instituciones políticas para mejorar procesos de toma de decisiones, demostrando compromiso ético y vocación de servicio público.
\end{itemize}

\subsection{Competencias Específicas}
Al término del curso, el estudiante tendrá la capacidad para:
\begin{enumerate}
    \item Identificar y justificar la estrategia metodológica óptima (RDiT, EHA, fsQCA, Text-as-Data o Teoría de Juegos) frente a una pregunta de investigación empírica doctoral, apoyándose en la lectura de referencia de esa ruta como modelo de redacción.
    \item Aplicar la lógica de identificación, especificación y, cuando sea pertinente a la ruta elegida, evaluación de robustez, replicando su estructura argumentativa en un proyecto propio.
    \item Implementar pipelines computacionales reproducibles en Python para la curaduría, estimación, diagnóstico de robustez y visualización de datos políticos y modelos estratégicos.
    \item Desarrollar un ensayo empírico cuantitativo (\textit{working paper}) para su tema de tesis, integrando la formulación del problema, la estimación del modelo y el análisis crítico de la evidencia, mediante avances progresivos revisados en asesoría.
\end{enumerate}

\subsection{Nota sobre la organización del curso}
El curso tiene 16 estudiantes divididos en dos grupos de 8 (Grupo A y Grupo B), que se reúnen de forma virtual los sábados de manera alterna: cada grupo asiste a un sábado de cada dos. Cada estudiante elige su ruta metodológica y su lectura de referencia al término de la sesión presencial magistral de la Semana 1. A partir de esa elección, cada estudiante trabaja de manera \textbf{autónoma entre asesorías}, siguiendo la Guía integrada de autoestudio correspondiente a su ruta, y presenta en cada sábado el entregable específico definido por la rúbrica de esa asesoría (Datos, Resultado o Interpretación, según corresponda): no son sesiones de contenido libre, sino \textbf{instancias evaluadas} con un entregable predefinido, cuyo avance recibe una calificación que forma parte de la nota final. La Semana 2 no tiene sesión presencial: queda para que cada estudiante, ya con su ruta y lectura de referencia elegidas en la Semana 1, trabaje de forma autónoma siguiendo la guía, avanzando hacia el entregable de la Asesoría 1.

\noindent
\textbf{Asistencia a asesorías y su peso en la nota.} La asistencia a las sesiones de asesoría es \textbf{obligatoria}. La nota de la \textbf{Entrega de Versión Final} (semana 9 para el Grupo A / semana 10 para el Grupo B) se compone de dos partes, hasta un máximo de \textbf{20 puntos}:

\begin{enumerate}
    \item \textbf{Hasta 16 puntos}, según el avance evaluado en las tres asesorías previas de cada grupo, ponderadas 40\%-30\%-30\% entre sí:
    \begin{itemize}
        \item \textbf{Asesoría 1 -- Datos} (semana 3 -- Grupo A / semana 4 -- Grupo B): \textbf{hasta 6.4 puntos}. Entregables: (i) documentos digitalizados (Word, PDF, Excel) a utilizar; (ii) archivo listo para leerse en Python (Excel, CSV o TXT), apoyándose en los notebooks de pre-procesamiento de la guía; (iii) archivo resultado listo para aplicar la técnica. \textit{Criterio:} se evalúa según el nivel de avance alcanzado en cada uno de estos tres componentes.
        \item \textbf{Asesoría 2 -- Resultado} (semana 5 -- Grupo A / semana 6 -- Grupo B): \textbf{hasta 4.8 puntos}. Entregables: (i) aplicación de la técnica; (ii) cuadros de resultados. \textit{Criterio:} se evalúa según el nivel de avance alcanzado en cada uno de estos dos componentes.
        \item \textbf{Asesoría 3 -- Interpretación} (semana 7 -- Grupo A / semana 8 -- Grupo B): \textbf{hasta 4.8 puntos}. Entregable: ensayo interpretativo según la(s) teoría(s) elegida(s), a partir de lo aprendido de los resultados. \textit{Criterio:} se evalúa según el nivel de avance del ensayo.
    \end{itemize}
    \item \textbf{Hasta 4 puntos adicionales}, evaluados en la propia Entrega de Versión Final (7 o 14 de noviembre), según la buena organización del documento presentado --- es decir, el cumplimiento de la Estructura Mínima del Working Paper (ver siguiente subsección).
\end{enumerate}

\noindent
\textbf{Inasistencia como criterio independiente de la nota.} Más allá del cálculo anterior, \textbf{no haber trabajado al inicio del curso es, por sí mismo, causal de desaprobación}: faltar a las \textbf{dos primeras asesorías} del propio grupo (semanas 3 y 5 para el Grupo A; semanas 4 y 6 para el Grupo B) antes del 7 de noviembre produce la \textbf{desaprobación automática del curso}, independientemente de la nota que se hubiera obtenido en cualquier otro componente. Esta regla no se deriva del cálculo ponderado de puntos: es un criterio administrativo aparte, que opera aun si --hipotéticamente-- el estudiante lograra puntaje en las demás instancias. Faltar a cualquier otra asesoría individual, sin activar este criterio, sí reduce el puntaje del componente correspondiente y eleva la probabilidad de desaprobar por la vía del cálculo ponderado.

\subsection{Estructura del Working Paper}
El documento que cada estudiante presenta en la \textbf{Entrega de Versión Final} (semana 9 -- Grupo A / semana 10 -- Grupo B) y, de corresponder, en la \textbf{reentrega} de diciembre, debe seguir como \textbf{estructura mínima} las siguientes secciones:

\begin{enumerate}
    \item \textbf{Presentación:} por qué este trabajo es relevante.
    \item \textbf{Pregunta de investigación:} qué busca responder el trabajo.
    \item \textbf{Marco teórico:} qué teoría o teorías busca verificar o mejorar.
    \item \textbf{Hipótesis o proposiciones teóricas:} lo que el marco teórico propone como respuesta a la pregunta de investigación (hipótesis, configuración causal esperada o equilibrio predicho, según la ruta elegida).
    \item \textbf{Diseño:} pasos a seguir para responder la pregunta; datos; unidad de análisis y de observación; técnica a emplear.
    \item \textbf{Resultados:} presentación de los resultados obtenidos, sean o no consistentes con la hipótesis planteada. Cuando sea apropiado para la ruta metodológica, se puede incluir un análisis de robustez que matice el resultado principal.
    \item \textbf{Interpretación:} qué dicen los resultados obtenidos.
    \item \textbf{Conclusiones.}
\end{enumerate}
Esta estructura es el formato mínimo exigido para dar cuerpo a lo que el estudiante reporta en la Entrega de Versión Final y, de ser necesario, en la reentrega; no reemplaza el contenido específico que cada ruta metodológica requiera.

\section{III. UNIDADES DE APRENDIZAJE}

\small
\SetTblrTemplate{caption}{empty}
\SetTblrTemplate{capcont}{empty}
\SetTblrTemplate{conthead}{empty}
\SetTblrTemplate{contfoot}{empty}
\begin{longtblr}{colspec={|Q[l,wd=1.0cm]|Q[l,wd=1.6cm]|Q[l,wd=3.6cm]|Q[l,wd=5.5cm]|Q[l,wd=2.4cm]|},rowhead=1}
\hline
\textbf{Sem.} & \textbf{Fecha} & \textbf{Unidad} & \textbf{Actividad} & \textbf{Material / Lectura} \\ \hline

01 & 12 set. & 
\textbf{Masterclass} (Grupos A y B) & 
Presentación de las cinco rutas metodológicas con datos ficticios y código Python end-to-end, siguiendo la guía del curso. Cada estudiante elige su ruta y su lectura de referencia. & 
Magallanes Reyes (2026), \textit{Métodos Cuantitativos de Análisis Político}. \\ \hline

02 & 19 set. & 
\textbf{Semana libre} (sin sesión) & 
Preparación autónoma, apoyándose en el capítulo de la guía y la lectura de referencia elegida. & 
Guía (capítulo de la ruta elegida); lectura de referencia elegida. \\ \hline

03 & 26 set. & \textbf{Asesoría 1 -- Datos} (Grupo A) & Entregable: documentos digitalizados, archivo legible en Python y archivo listo para la técnica (ver rúbrica en Sección VI). Asesoría obligatoria y evaluada: \textbf{hasta 6.4 puntos} de los 20 de la Entrega de Versión Final (semana 9). Faltar a esta sesión y a la Asesoría 2 produce desaprobación automática (criterio aparte de la nota).  & Magallanes Reyes (2026), \textit{Métodos Cuantitativos de Análisis Político}. \\ \hline
04 & 3 oct.  & \textbf{Asesoría 1 -- Datos} (Grupo B) & Entregable: documentos digitalizados, archivo legible en Python y archivo listo para la técnica (ver rúbrica en Sección VI). Asesoría obligatoria y evaluada: \textbf{hasta 6.4 puntos} de los 20 de la Entrega de Versión Final (semana 10). Faltar a esta sesión y a la Asesoría 2 produce desaprobación automática (criterio aparte de la nota).  & Magallanes Reyes (2026), \textit{Métodos Cuantitativos de Análisis Político}. \\ \hline
05 & 10 oct. & \textbf{Asesoría 2 -- Resultado} (Grupo A) & Entregable: aplicación de la técnica y cuadros de resultados (ver rúbrica en Sección VI). Asesoría obligatoria y evaluada: \textbf{hasta 4.8 puntos} de los 20 de la Entrega de Versión Final (semana 9). Faltar a esta sesión y a la Asesoría 1 produce desaprobación automática (criterio aparte de la nota).  & Magallanes Reyes (2026), \textit{Métodos Cuantitativos de Análisis Político}. \\ \hline
06 & 17 oct. & \textbf{Asesoría 2 -- Resultado} (Grupo B) & Entregable: aplicación de la técnica y cuadros de resultados (ver rúbrica en Sección VI). Asesoría obligatoria y evaluada: \textbf{hasta 4.8 puntos} de los 20 de la Entrega de Versión Final (semana 10). Faltar a esta sesión y a la Asesoría 1 produce desaprobación automática (criterio aparte de la nota).  & Magallanes Reyes (2026), \textit{Métodos Cuantitativos de Análisis Político}. \\ \hline
07 & 24 oct. & \textbf{Asesoría 3 -- Interpretación} (Grupo A) & Entregable: ensayo interpretativo según la(s) teoría(s) elegida(s), a partir de lo aprendido de los resultados (ver rúbrica en Sección VI). Asesoría obligatoria y evaluada: \textbf{hasta 4.8 puntos} de los 20 de la Entrega de Versión Final (semana 9).  & Magallanes Reyes (2026), \textit{Métodos Cuantitativos de Análisis Político}. \\ \hline
08 & 31 oct. & \textbf{Asesoría 3 -- Interpretación} (Grupo B) & Entregable: ensayo interpretativo según la(s) teoría(s) elegida(s), a partir de lo aprendido de los resultados (ver rúbrica en Sección VI). Asesoría obligatoria y evaluada: \textbf{hasta 4.8 puntos} de los 20 de la Entrega de Versión Final (semana 10).  & Magallanes Reyes (2026), \textit{Métodos Cuantitativos de Análisis Político}. \\ \hline

09 & 7 nov. & 
\textbf{Entrega de Versión Final} (evaluación, Grupo A) & 
\textbf{[Entrega del Working Paper --- versión final]}, Grupo A. Nota sobre 20: hasta 16 puntos por las Asesorías 1, 2 y 3, más hasta 4 puntos por la organización del documento.  & Magallanes Reyes (2026), \textit{Métodos Cuantitativos de Análisis Político}. \\ \hline

10 & 14 nov. & 
\textbf{Entrega de Versión Final} (evaluación, Grupo B) & 
\textbf{[Entrega del Working Paper --- versión final]}, Grupo B. Nota sobre 20: hasta 16 puntos por las Asesorías 1, 2 y 3, más hasta 4 puntos por la organización del documento.  & Magallanes Reyes (2026), \textit{Métodos Cuantitativos de Análisis Político}. \\ \hline

11 & 21 nov. & \textbf{Asesoría Adicional / Cierre} (Grupo A) & Para nota $<14$: asesoría adicional obligatoria, previa a la reentrega. Para nota $\geq 14$: preparación de la sustentación.  & Magallanes Reyes (2026), \textit{Métodos Cuantitativos de Análisis Político}. \\ \hline
12 & 28 nov. & \textbf{Asesoría Adicional / Cierre} (Grupo B) & Para nota $<14$: asesoría adicional obligatoria, previa a la reentrega. Para nota $\geq 14$: preparación de la sustentación.  & Magallanes Reyes (2026), \textit{Métodos Cuantitativos de Análisis Político}. \\ \hline

13 & 5 dic. & 
\textbf{Reentrega} (evaluación, Grupo A) & 
\textbf{[Reentrega de la versión observada]}, solo para quienes obtuvieron nota $<14$ en la Entrega de Versión Final. La nota puede mejorar según el avance mostrado en la Asesoría Adicional.  & Magallanes Reyes (2026), \textit{Métodos Cuantitativos de Análisis Político}. \\ \hline

14 & 12 dic. & 
\textbf{Reentrega} (evaluación, Grupo B) & 
\textbf{[Reentrega de la versión observada]}, solo para quienes obtuvieron nota $<14$ en la Entrega de Versión Final. La nota puede mejorar según el avance mostrado en la Asesoría Adicional.  & Magallanes Reyes (2026), \textit{Métodos Cuantitativos de Análisis Político}. \\ \hline

15 & 19 dic. & 
\textbf{Exposición} (nota final $\geq 14$) & 
Sustentación pública del working paper (15 min + 10 min de debate), para quienes alcanzaron nota final $\geq 14$ (en la Entrega de Versión Final de noviembre o tras la reentrega de diciembre). & 
Working papers de los estudiantes. \\ \hline

16 & 26 dic. & 
\textbf{Examen Recuperatorio} (nota final $< 14$ tras reentrega) & 
Examen escrito sobre \textbf{toda la Guía integrada de autoestudio} del curso. Solo pueden rendirlo quienes asistieron a todas sus asesorías y presentaron todas sus entregas, aun incompletas. Si se aprueba con nota $>14$, la nota final del curso queda topada en \textbf{14}. & 
Guía integrada de autoestudio (todos los capítulos). \\ \hline
\end{longtblr}
\normalsize

\noindent
\textit{La calificación final y el envío de notas a la Unidad de Posgrado se realiza como trámite administrativo posterior al 26 de diciembre, fuera de las sesiones del curso.}

\section{IV. MEDIOS Y MATERIALES}
Para el desarrollo de la asignatura se empleará la plataforma virtual Moodle como repositorio central de recursos y gestor de actividades, complementada con Google Meet para sesiones sincrónicas de asesoría. El soporte computacional se basará en Python (Google Colab; librerías \texttt{statsmodels}, \texttt{lifelines}, \texttt{setqca}, \texttt{sentence-transformers}) y GitHub para garantizar estándares internacionales de replicabilidad.

\section{V. ESTRATEGIAS METODOLÓGICAS}
Las estrategias metodológicas que se van a utilizar en la asignatura son:
\begin{itemize}
    \item \textbf{Masterclass y Modelado Demostrativo (Semana 1):} Presentación intensiva de las cinco rutas analíticas con código y datos ficticios, siguiendo la guía del curso.
    \item \textbf{Autoestudio guiado:} Recorrido individual del capítulo de la guía correspondiente a la ruta elegida, con notebooks reproducibles de preparación y análisis.
    \item \textbf{Asesoría obligatoria y evaluada por grupos, con entregable definido por rúbrica:} Sesiones de asesoría virtual en dos grupos de 8 estudiantes; entre cada asesoría, el estudiante trabaja de manera autónoma siguiendo la guía, avanzando hacia el entregable específico de la siguiente sesión (Datos, Resultado o Interpretación). Las tres primeras asesorías de cada grupo aportan hasta 16 de los 20 puntos de la Entrega de Versión Final; los 4 puntos restantes se evalúan según la organización del documento presentado.
    \item \textbf{Elaboración progresiva de un Working Paper:} Una única entrega de versión final, con posibilidad de reentrega y mejora de nota para quien no alcance la nota mínima.
    \item \textbf{Cierre diferenciado por desempeño:} Sustentación pública para quienes aprueban, o examen recuperatorio sobre toda la guía del curso para quienes no alcanzan la nota mínima tras la reentrega.
\end{itemize}

\section{VI. EVALUACIÓN DE LOS APRENDIZAJES}

La nota final del curso corresponde a la nota de la \textbf{Entrega de Versión Final} del working paper (semana 9 -- Grupo A / semana 10 -- Grupo B), sobre un máximo de \textbf{20 puntos}:

\begin{table}[ht!]
\centering
\begin{tabularx}{\textwidth}{|X|X|X|c|}
\hline
\textbf{COMPONENTE} & \textbf{ENTREGABLE} & \textbf{CRITERIO} & \textbf{PUNTAJE} \\ \hline
Asesoría 1 -- Datos (semana 3 -- Grupo A / semana 4 -- Grupo B) & Documentos digitalizados; archivo legible en Python; archivo listo para la técnica & Nivel de avance en cada uno de los tres componentes & \textbf{hasta 6.4} \\ \hline
Asesoría 2 -- Resultado (semana 5 -- Grupo A / semana 6 -- Grupo B) & Aplicación de la técnica; cuadros de resultados & Nivel de avance en cada uno de los dos componentes & \textbf{hasta 4.8} \\ \hline
Asesoría 3 -- Interpretación (semana 7 -- Grupo A / semana 8 -- Grupo B) & Ensayo interpretativo según la(s) teoría(s) elegida(s) & Nivel de avance del ensayo & \textbf{hasta 4.8} \\ \hline
Organización del documento (Entrega, 7/14 nov.) & Cumplimiento de la Estructura Mínima del Working Paper & Buena organización del documento presentado & \textbf{hasta 4} \\ \hline
\multicolumn{3}{|r|}{\textbf{TOTAL --- Nota de la Entrega de Versión Final}} & \textbf{20} \\ \hline
\end{tabularx}
\end{table}

\noindent
El documento entregado debe seguir la \textbf{estructura mínima del Working Paper} definida en la Sección II.

\noindent
\textbf{Todos los estudiantes deben presentar la Entrega de Versión Final} de noviembre; no existe una vía alternativa que exima de esta obligación.

\noindent
\textbf{Mejora de la nota (semanas 11 a 14).} El estudiante cuya nota de la Entrega de Versión Final sea \textbf{menor a 14} recibe una \textbf{asesoría adicional} obligatoria (semana 11 -- Grupo A / semana 12 -- Grupo B) y debe presentar una \textbf{reentrega de la versión observada} (semana 13 -- Grupo A / semana 14 -- Grupo B). La nota puede mejorar en función del avance mostrado en la asesoría adicional, verificado mediante dicha reentrega. Para los estudiantes que ya alcanzaron nota $\geq 14$ en noviembre, estas semanas no implican una nueva evaluación obligatoria: la asesoría se dedica a la preparación de la sustentación.

\noindent
\textbf{Sustentación (Semana 15).} Los estudiantes cuya nota final sea \textbf{igual o mayor a 14} --- ya sea en la Entrega de Versión Final de noviembre o tras la mejora registrada en la reentrega de diciembre --- sustentan públicamente su working paper. Esta sustentación no añade ni resta puntaje: formaliza el cierre del curso para quienes ya aprobaron.

\noindent
\textbf{Examen recuperatorio (Semana 16).} Los estudiantes cuya nota permanezca \textbf{menor a 14} después de la reentrega de diciembre rinden un examen escrito sobre \textbf{toda la Guía integrada de autoestudio del curso} (los capítulos correspondientes a las cinco rutas metodológicas), independientemente de la ruta elegida para su propio proyecto. El examen recuperatorio \textbf{no tiene un porcentaje propio} dentro de la nota final: es un mecanismo de segunda oportunidad para alcanzar la nota mínima aprobatoria, no un componente adicional de la evaluación regular. \textbf{Tope de nota:} si el estudiante aprueba el examen recuperatorio con una nota mayor a 14, su nota final del curso se registra como \textbf{14 (nota máxima por esta vía)}; el examen recuperatorio permite alcanzar la condición de aprobado, pero no permite obtener una nota final superior a la mínima aprobatoria.

\noindent
\textbf{Restricción de acceso al examen recuperatorio.} Solo puede rendir el examen recuperatorio quien (i) asistió a \textbf{todas} las sesiones de asesoría obligatorias de su grupo y (ii) presentó \textbf{todas} sus entregas (la Entrega de Versión Final de noviembre y, de corresponder, la reentrega de diciembre), aun cuando estén incompletas. El estudiante que no cumpla ambas condiciones \textbf{no accede} al examen recuperatorio, y su nota --- la de la reentrega de diciembre si la hubo, o la de la Entrega de Versión Final de noviembre en caso contrario --- es definitiva.

\noindent
\textbf{Desaprobación automática por inasistencia (criterio independiente de la nota).} Más allá del cálculo de puntos descrito arriba, no haber trabajado desde el inicio del curso es, por sí mismo, causal de desaprobación: faltar a las \textbf{dos primeras asesorías} del propio grupo (semanas 3 y 5 para el Grupo A; semanas 4 y 6 para el Grupo B) antes del 7 de noviembre produce la \textbf{desaprobación automática} del curso, independientemente del puntaje que se hubiera obtenido en cualquier otro componente. Este criterio es administrativo, no se deriva del cálculo ponderado de puntos, y prevalece sobre él. Faltar a cualquier otra asesoría individual, sin activar este criterio, sí reduce el puntaje del componente correspondiente y eleva la probabilidad de desaprobar por la vía del cálculo de puntos.

De acuerdo con el artículo 150 del Reglamento General de Estudios de Posgrado, la calificación de las asignaturas cursadas por los estudiantes se rige por la escala vigesimal. La mínima nota aprobatoria por asignatura es catorce (14). El promedio ponderado de las asignaturas, al concluir los estudios, debe ser mínimo catorce (14).

\section{VII. BIBLIOGRAFÍA DEL CURSO}

\subsection{Guía del curso}
\begin{itemize}
    \item Magallanes Reyes, J. M. (2026). \textit{Métodos Cuantitativos de Análisis Político: Guía integrada de autoestudio}. Universidad Nacional Mayor de San Marcos, Doctorado en Ciencia Política. \textbf{Base íntegra del examen recuperatorio (Sección VI).}
\end{itemize}

\subsection{Lecturas de Referencia por Ruta Metodológica}
Cada estudiante elige, de esta lista, la lectura correspondiente a su ruta metodológica, como modelo de redacción para su propio working paper. Estas lecturas no forman parte del contenido del examen recuperatorio, el cual se basa exclusivamente en la Guía integrada de autoestudio del curso.

\begin{itemize}
    \item \textbf{[Ruta 1 -- RDiT]:} Reny, T. T., \& Newman, B. J. (2021). The Opinion-Mobilizing Effect of Social Protest against Police Violence: Evidence from the 2020 George Floyd Protests. \textit{American Political Science Review}, 115(4), 1499--1507.
    \item \textbf{[Ruta 2 -- EHA]:} Feliú Ribeiro, P., \& López Burian, C. (2023). Factores domésticos en la supervivencia de los ministros de Relaciones Exteriores en América Latina (1945-2020). \textit{Estado, Gobierno y Gestión Pública}, 21(40).
    \item \textbf{[Ruta 3 -- fsQCA]:} Vis, B. (2012). The Comparative Advantages of fsQCA and Regression Analysis for Moderately Large-N Analyses. \textit{Sociological Methods \& Research}, 41(1), 168--198.
    \item \textbf{[Ruta 4 -- Text-as-Data]:} Ceron, T., Blokker, N., \& Padó, S. (2022). Optimizing text representations to capture (dis)similarity between political parties. \textit{Proceedings of the 26th Conference on Computational Natural Language Learning (CoNLL)}, 325--338.
    \item \textbf{[Ruta 5 -- Teoría de Juegos]:} Clinton, R. L. (1994). Game Theory, Legal History, and the Origins of Judicial Review: A Revisionist Analysis of Marbury v. Madison. \textit{American Journal of Political Science}, 38(2), 285--302.
\end{itemize}

\vspace{2em}
\begin{flushright}
    Ciudad Universitaria, 27 de agosto de 2026
\end{flushright}

\end{document}
